\documentclass[12pt]{article}
\usepackage[margin=0.3in,top=0.25in,bottom=0.4in]{geometry}
\makeatletter
\renewcommand{\@maketitle}{\newpage\null\vskip -3.5em\begin{center}{\LARGE \@title \par}\end{center}\par\vskip 0.2em}
\makeatother
\usepackage{amsmath,amssymb}
\usepackage{booktabs}
\usepackage{hyperref}
\setlength{\parskip}{0.6em}
\setlength{\parindent}{0pt}
\newcommand{\xpo}{x_{\mathbb{P}}}
\newcommand{\dd}{\mathrm{d}}

\title{The quark DTMD in the diffractive part}
\author{}
\date{}

\begin{document}
\maketitle

\section*{In short}

The diffractive (TMD) part of the code now has two terms, the gluon DTMD (Sec.~3 of Ref.~[3],
Caucal, Gimeno-Estivill, Iancu, Lappi, Salazar, arXiv:2606.04169) and the quark DTMD (Sec.~4 of
Ref.~[3]), both taken at photon virtuality $Q^2=0$. They share the prefactor that the plotting scripts
already apply,
\[
  \mathcal{P}=\frac{\alpha_{\mathrm{em}}\,\alpha_s\,e_c^2\,(N_c^2-1)\,S_\perp}{8\pi^4}\,,
\]
so only the C++ integrand changes. (In the scripts $S_\perp$ comes from the $b_d$ integral, or from
$\sigma_0/2$ for the proton.) The environment variable \texttt{DTMD} selects \texttt{gluon},
\texttt{quark} or \texttt{both} (default).

\section*{1. Which configuration}

In the quark DTMD the semi-hard parton (transverse momentum $K_\perp\sim Q_s$) is a quark, and the
hard dijet ($P_\perp\gg Q_s$) is the antiquark and the gluon. Ref.~[3] computes the case of a soft quark;
by charge conjugation the soft antiquark gives the same with the roles exchanged. For $D^0$ production
the code uses the configuration in which the \textbf{measured charm is hard}:
\begin{itemize}
\item the charm $c$ has transverse momentum $p\simeq P_\perp$ and longitudinal fraction $z=p^+/q^+$
  (Ref.~[3]: the hard parton 2, $z_2=z$);
\item the gluon is hard, with fraction $\bar z=1-z$ (Ref.~[3]: $z_3=\bar z$);
\item the $\bar c$ is semi-hard, with momentum $K_\perp$ and $z_1\simeq0$.
\end{itemize}
This configuration is counted once. That is right for $D^0$ alone and for $(D^0+\bar D^0)/2$. The
semi-hard charm itself (at $K_\perp\sim Q_s$) is not counted as a source of $D^0$.

\section*{2. The formula of Ref.~[3]}

Eqs.~(4.52)--(4.54) of Ref.~[3]:
\[
  \frac{\dd\sigma}{\dd z_2\,\dd z_3\,\dd^2K\,\dd^2P\,\dd\ln(1/x)}
  =\frac{4\pi^2\alpha_{\mathrm{em}}}{Q^2}\,e_f^2\,\mathcal{H}(z_2,z_3,P_\perp,Q)\,
  \frac{\dd\,x q_{\mathbb{P}}(x,\xpo,K_\perp,m)}{\dd^2K}\,,
\]
\[
  \mathcal{H}=\delta(1-z_2-z_3)\,\frac{\alpha_s C_F}{2\pi^2}\,\frac{1}{z_3}\,
  \frac{\hat Q^2\big[(P_\perp^2+\hat Q^2)^2+P_\perp^4z_3^2+\hat Q^4z_2^2\big]}{P_\perp^2\,(P_\perp^2+\hat Q^2)^3}\,,
  \qquad \hat Q^2=z_2z_3Q^2\,,
\]
\[
  \frac{\dd\,x q_{\mathbb{P}}}{\dd^2K}=\frac{S_\perp N_c}{8\pi^4}\,\frac{x}{1-x}\,
  \frac{K_\perp^2+m^2}{xK_\perp^2+m^2}
  \left(\mathcal{Q}_1^2+\frac{(1-x)\,m^2}{xK_\perp^2+m^2}\,\mathcal{Q}_0^2\right).
\]
With Eqs.~(4.21), (4.23) and (4.50),
\[
  \mathcal{Q}_n(\omega,K_\perp)=\omega^2\int_0^\infty\dd R\,R\,J_n(K_\perp R)\,K_n(\omega R)\,
  \mathcal{N}(R,\xpo)\quad(n=0,1)\,,\qquad
  \omega^2=\frac{m^2+xK_\perp^2}{1-x}\,.
\]
$\mathcal{N}$ is the \textbf{quark} (fundamental) dipole amplitude: the semi-hard quark and the compact
$\bar qg$ pair scatter as a $q\bar q$ dipole of size $R$. (The gluon DTMD has instead
$\mathcal{N}_g=2\mathcal{N}-\mathcal{N}^2$.) In the code: \texttt{dipole->N(r, x)}, at
$\min(\xpo,0.01)$ as in the other parts.

\section*{3. The limit $Q^2\to0$}

The $1/Q^2$ in front cancels against $\hat Q^2=z_2z_3Q^2$ in $\mathcal{H}$, and for $\hat Q\to0$ the
bracket is $P_\perp^4(1+z_3^2)$:
\[
  \frac{4\pi^2\alpha_{\mathrm{em}}}{Q^2}\,e_f^2\,\mathcal{H}\;\xrightarrow{\;Q^2\to0\;}\;
  4\pi^2\alpha_{\mathrm{em}}e_f^2\,\frac{\alpha_sC_F}{2\pi^2}\,\frac{z_2z_3}{z_3}\,
  \frac{P_\perp^4(1+z_3^2)}{P_\perp^2P_\perp^6}
  =2\,\alpha_{\mathrm{em}}\alpha_s e_f^2\,C_F\,\frac{z\,(1+\bar z^2)}{p^4}\,.
\]
The limit is finite: in photoproduction the hard scale is $P_\perp$, through $1/P_\perp^4$.

The longitudinal fraction is traded for $x$ with Eq.~(4.48), $x=x_{\bar qg}/\xpo$. At $Q^2=0$, and with
the charm (mass $m$, fraction $z$) and the gluon ($\bar z$) as the hard pair,
\[
  x=\frac{M^2}{\xpo W^2}\,,\qquad M^2=\frac{p^2}{z\bar z}+\frac{m^2}{z}\,,
\]
the invariant mass of the charm--gluon pair. $\omega^2$ above does not depend on $Q$.

\section*{4. The same prefactor as the gluon DTMD}

Multiplying the hard factor of Sec.~3 with the quark DTMD,
\[
  2\,\alpha_{\mathrm{em}}\alpha_s e_f^2\,C_F\,\frac{z(1+\bar z^2)}{p^4}\cdot\frac{S_\perp N_c}{8\pi^4}\,F_q
  =\frac{\alpha_{\mathrm{em}}\alpha_s e_f^2\,S_\perp}{8\pi^4}\,(2C_FN_c)\,\frac{z(1+\bar z^2)}{p^4}\,F_q
  =\mathcal{P}\,\frac{z(1+\bar z^2)}{p^4}\,F_q\,,
\]
since $2C_FN_c=N_c^2-1$, with
\[
  F_q=\frac{x}{1-x}\,\frac{K_\perp^2+m^2}{xK_\perp^2+m^2}
  \left(\mathcal{Q}_1^2+\frac{(1-x)\,m^2}{xK_\perp^2+m^2}\,\mathcal{Q}_0^2\right).
\]
For comparison, the gluon DTMD (Eqs.~(3.38), (3.39), (3.43)--(3.45) of Ref.~[3], at $Q^2=0$) is
\[
  \mathcal{P}\,\frac{(z^2+\bar z^2)(p^4+m^4)+2m^2p^2}{(p^2+m^2)^4}\,\frac{G^2}{1-x}\,,\qquad
  G=\Omega^2\!\int_0^\infty\!\dd R\,R\,J_2(K_\perp R)K_2(\Omega R)\,\mathcal{N}_g(R,\xpo)\,,\quad
  \Omega^2=\frac{x}{1-x}K_\perp^2\,,
\]
with $x=M_{c\bar c}^2/(\xpo W^2)$ and $M_{c\bar c}^2=(p^2+m^2)/(z\bar z)$.

\section*{5. What the code computes}

\texttt{quark\_dtmd\_kernel} in \texttt{src/integrands.cpp} returns, per $\dd\ln(1/x)=\dd\ln\xpo$ and
without $\mathcal{P}$,
\[
  2\pi K_\perp\;z\;\frac{z(1+\bar z^2)}{p^4}\;F_q\,,
\]
where
\begin{itemize}
\item $2\pi K_\perp$ is $\dd^2K=2\pi K_\perp\dd K_\perp$ (the angle is trivial);
\item the first $z$ is the Jacobian $\dd z/\dd y=z$ at fixed $q^+$, as in the gluon part;
\item $\mathcal{Q}_n^2$ is computed as in the gluon part, with the two radial integrals as the two
  VEGAS variables $r_1,r_2$:
  $\mathcal{Q}_n^2\to\omega^4\,r_1r_2\,J_n(K_\perp r_1)K_n(\omega r_1)\mathcal{N}(r_1)\,
  J_n(K_\perp r_2)K_n(\omega r_2)\mathcal{N}(r_2)$.
\end{itemize}
Everything else is shared with the gluon part: the photon flux, the fragmentation $c\to D^0$, the
logarithmic sampling of $K_\perp$ and $\xpo$, the cut $K_\perp\leq c\,p$ (\texttt{K\_CUT\_FACTOR}) and
$\alpha_s(m_T)$ in $\mathcal{P}$.

\section*{6. Caveats}

\begin{itemize}
\item \textbf{Massless hard factor.} In Ref.~[3] the quark mass is a semi-hard scale in this
  configuration: it enters $x$, $\omega$ and $F_q$, but not $\mathcal{H}$. The quark DTMD term therefore
  needs $p\gg m_c$ as well as $p\gg Q_s$, while the gluon DTMD term has $m_c$ in its hard factor.
\item \textbf{Counting of $D^0$} as in Sec.~1 (one configuration, measured charm hard).
\item \textbf{Old results} have the gluon DTMD only. New output files have a header line
  \texttt{DTMD : both/gluon/quark}.
\end{itemize}

\end{document}
